
\begin{figure}[H]
\centering
\scriptsize
\setlength{\tabcolsep}{2pt}
\renewcommand{\arraystretch}{1.25}
\begin{tabular}{L{0.18\textwidth}L{0.15\textwidth}L{0.20\textwidth}L{0.15\textwidth}L{0.20\textwidth}}
\toprule
\textbf{Observable} & \textbf{left: $\rho_O<0$} & \textbf{diagonal: $\rho_O\simeq0$} & \textbf{right: $\rho_O>0$} & \textbf{unmeasured / status}\\
\midrule
$m/q$ & -- & \cellcolor{green!45}$\anti{H}$ / $\bar p$ compatible & -- & \cellcolor{orange!35}$\anti{He}$: K; $Z\geq3$: ?\\
$1S$--$2S$ & -- & \cellcolor{green!45}$\anti{H}[1S\text{--}2S:0]$ & -- & \cellcolor{gray!8}$Z\geq2$: ?\\
$g$ & -- & \cellcolor{green!45}$\anti{H}[g:0]$ & -- & \cellcolor{gray!8}$Z\geq2$: ?\\
Chemistry & -- & -- & -- & \cellcolor{yellow!35}theoretically prepared / ?\\
Cosmic search & -- & -- & -- & \cellcolor{yellow!35}search domain, not anti-atom chemistry\\
\bottomrule
\end{tabular}

\vspace{0.45em}
\scriptsize
The diagonal means CPT compatibility within uncertainty. Left/right entries are used only when a significant negative or positive residual is measured. Missing data remain marked as ? or K and are not interpreted as deviations.
\caption{Variant C: diagonal and residual scheme for possible matter-antimatter asymmetries. Residual entries are tied only to named measured observables or explicitly hypothetical scenarios; unknown anti-systems stay separately marked as unmeasured.}
\label{fig:asymmetry-map}
\end{figure}
